\chapter[Published optical-trap measurements]{An exploratory analysis of published optical-trap measurements}
\label{ch:ci-empirical}

\section{The experimental question}

Published measurements allow a practical test of the distinction between
observing correlations and having enough calibration to certify dissipation.
We analyze three records from Bryan and colleagues' optical-trap study
\cite{ciBryan2022,ciBryanData2022}. Two records measure the same particle
simultaneously through detectors with different noise levels. The third is
a separate harmonic-well record. The source describes harmonic confinement
without feedback, while feedback creates other potentials in the study.

The useful result is a detector-calibration diagnostic. The simultaneous
records support a conditional decomposition of position and measurement
variance, and show why the lower-noise channel cannot simply be treated as
exact position. They also provide an empirical example in which small
correlation-moment departures do not justify a thermodynamic certificate.
No entropy-production confidence bound is issued for these data.

The repository is distributed under CC-BY-4.0 and credits J.~S.~Bryan~IV,
P.~Basak, J.~Bechhoefer and S.~Press\'e. We preserve the original files and
verify their published checksums. The source filename
\texttt{harmwell\_n1\_groundtruth.h5} denotes the simultaneous lower-noise
measurement. It does not provide exact latent positions. Serialized force
functions in that file are also distinct from measured force trajectories;
the loader never executes them.

\section{Fixed processing and its limits}

Each file contains $49\,999$ positions. The authors' code documents a nominal
sampling interval of $10\,\mu\mathrm{s}$. Individual timestamps are absent.
We use the first $10\,000$ points for centering, scale selection and a simple
covariance fit, and the remaining $39\,999$ for evaluation. The fixed lag menu,
in sample intervals, is
\[
  0,1,2,4,8,16,32,64,128,256.
\]
These choices preceded the new correlation calculations, but an earlier audit
had inspected the sources and whole-file descriptive statistics. This is an
exploratory reanalysis, not a prospectively blinded experiment.

For the simultaneous detectors, each training mean supplies its center and
twice the square root of the average training variance supplies a common
scale. For the separate record, we use its own training mean and twice its
training standard deviation. The bounded observation is
\[
  Z_i=\tanh\!\left(\frac{X_i-c_{\mathrm{train}}}
                                  {s_{\mathrm{train}}}\right).
\]
Neither center nor scale is refitted on the evaluation segment. The common
scale for the simultaneous detectors is $29.0285\,\mathrm{nm}$; the scale for
the separate record is $30.9707\,\mathrm{nm}$.

Every reported covariance uses the mean of its complete evaluation segment
and averages the $N-k$ available ordered pairs at lag $k$. We repeat the
calculation in eight contiguous blocks, using each block's own mean and
excluding pairs that cross a block boundary. These block ranges measure
sensitivity to the selected segment. They are not confidence intervals,
independent replications, or a calibration of hidden-state mixing. Mean
estimation itself can bias finite-trajectory covariances.

The evaluation span is nominally $0.39998\,\mathrm{s}$ per record, with an
additional $0.09999\,\mathrm{s}$ training span. The two simultaneous detectors
share the same physical exposure. Adding their sample counts would not double
that exposure. Overlapping lag pairs share observations and remain dependent;
the separate record's alignment relative to the first experiment is unknown.

\section{What two detectors can and cannot calibrate}

Consider the explicitly conditional model
\[
 X_t=S_t+\epsilon_t,\qquad Y_t=S_t+\eta_t,
\]
with unit gains and zero-mean errors independent of each other and of the
particle. Under these assumptions,
\begin{align}
 \operatorname{Var}(S)&=\operatorname{Cov}(X,Y),\\
 \operatorname{Var}(\epsilon)&=\operatorname{Var}(X)-\operatorname{Cov}(X,Y),\\
 \operatorname{Var}(\eta)&=\operatorname{Var}(Y)-\operatorname{Cov}(X,Y).
\end{align}
The evaluation data give the following plug-in decomposition. These values
are conditional model diagnostics, not independently calibrated variances.

\begin{center}
\begin{tabular}{lr}
\hline
Quantity & Estimate ($\mathrm{nm}^2$)\\
\hline
Variance of detector $X$ &234.411\\
Variance of lower-noise detector $Y$ &226.424\\
Same-time covariance &218.387\\
Candidate noise variance of $X$ &16.024\\
Candidate noise variance of $Y$ &8.037\\
Variance of $X-Y$ &24.061\\
\hline
\end{tabular}
\end{center}

Treating $Y$ as exact position would assign the entire difference variance
$24.061$ to the first detector. Under the stated independent-error model,
that is about $50.2\%$ larger than its candidate variance $16.024$. The
repository's distinction between a lower-noise reference and latent truth
therefore has a measurable consequence for this analysis. Correlated errors,
unequal gains or detector memory could change the interpretation; the
decomposition does not establish their absence.

Adding temporal whiteness to the model predicts zero positive-lag covariance
of $X-Y$ and equality of the two auto-covariances and two ordered
cross-covariances. The table reports the difference-record autocorrelation
and the difference of ordered cross-covariances at every declared positive
lag. The latter is $\widehat{\operatorname{Cov}}(X_t,Y_{t+k})-
\widehat{\operatorname{Cov}}(Y_t,X_{t+k})$.

\begin{center}
\begin{tabular}{rrrr}
\hline
Lag ($\mu\mathrm{s}$)&Difference correlation&Block range&Cross difference ($\mathrm{nm}^2$)\\
\hline
10   &0.0092&$[-0.0151,\phantom{-}0.0259]$& 0.195\\
20   &0.0214&$[-0.0064,\phantom{-}0.0541]$& 0.096\\
40   &0.0101&$[-0.0162,\phantom{-}0.0288]$& 0.160\\
80   &0.0150&$[-0.0063,\phantom{-}0.0319]$&-0.052\\
160  &0.0190&$[\phantom{-}0.0030,\phantom{-}0.0364]$&-0.168\\
320  &0.0076&$[-0.0147,\phantom{-}0.0270]$& 0.174\\
640  &0.0010&$[-0.0295,\phantom{-}0.0310]$& 0.146\\
1280 &0.0127&$[-0.0115,\phantom{-}0.0401]$&-0.155\\
2560 &-0.0009&$[-0.0105,\phantom{-}0.0078]$&0.410\\
\hline
\end{tabular}
\end{center}

The residual correlation at $160\,\mu\mathrm{s}$ is positive in all eight
blocks, so exact white measurement error should remain an assumption to
check. These descriptive signs do not constitute a statistical rejection:
block dependence and selection across the lag menu are uncalibrated. The
auto-covariance difference at the shortest lag is $3.881\,\mathrm{nm}^2$,
with a block range $[-0.692,7.712]$. Ordered cross-covariance differences
also have block ranges spanning zero at each tested positive lag.

There is a separate limitation for bounded observations. Even conditional
independence of the two channels would make their same-time product estimate
$\mathbb E[f_X(S)f_Y(S)]$, where $f_X,f_Y$ are their conditional means.
Different noise levels generally make these functions different after the
nonlinear transform. Consequently, this product is not automatically the
replicate calibration $A=\mathbb E[f_X(S)^2]$ used elsewhere in this work.
Matching detector responses is an additional requirement.

\section{A relaxation model and reversible moment checks}

A naive Ornstein--Uhlenbeck covariance model,
$C(t)=a\exp(-\lambda t)$, was fitted only to the symmetrized training
cross-covariances at lags $1,2,4,8,16$, using least squares on their logarithms.
It gives $a=198.967\,\mathrm{nm}^2$ and
$\lambda=1536.63\,\mathrm{s}^{-1}$, a relaxation time of
$0.6508\,\mathrm{ms}$. On evaluation data its shortest-lag prediction is
$195.933\,\mathrm{nm}^2$, compared with $215.798\,\mathrm{nm}^2$ observed.
The discrepancy is $9.10\%$ of the evaluation zero-lag cross-covariance.
This out-of-segment discrepancy matters if a fitted relaxation model is used
as calibration. It can reflect finite-trajectory fluctuation, drift or model
mismatch; these records alone do not distinguish them. In particular, the
fitted relaxation rate is not a bound on all hidden modes.

For each bounded detector, let $g(t)$ denote its centered positive-lag
covariance and $v$ its observed variance. A stationary reversible Markov
model with a memoryless detector gives a positive spectral mixture for $g$.
The detector's conditional variance contribution is at most $v$. The tested
necessary population inequalities are
\begin{align}
 g(t)&\geq0, & g(t)-g(2t)&\geq0,\\
 g(t)-g(2t)&\leq v/4, & g(2t)^3&\leq g(t)^2g(4t).
\end{align}
The last inequality uses the weights $2/3$ and $1/3$ in H\"older's
inequality, because $2t=(2/3)t+(1/3)4t$. It is not the equal-weight
midpoint inequality. The pure-exponential and spectral-mixture tests check
this convention explicitly.

All three full evaluation segments satisfy the first three checks at every
applicable declared lag. The smallest upper-drop margins are $0.00466$,
$0.00119$ and $0.00293$ for the first detector, reference and separate record,
respectively. Some empirical weighted log-convexity margins are negative,
but every such negative full-segment margin has a block range spanning zero.
Individual blocks also give negative values for some of the simpler checks.
These small-sample sensitivity results provide no certified violation of
the reversible class. Passing a finite menu would not establish equilibrium
or identify total dissipation either.

\section{What would enable a physical certificate}

The executable analysis returns \texttt{abstain} for thermodynamic inference.
Converting this workflow into a valid lower bound requires a physical and
statistical contract beyond a good fitted covariance:
\begin{enumerate}
\item Calibrate detector gain, temporal response and error uniformly over the
      working position range. For same-state replicate calibration, match
      conditional mean responses and justify conditional independence.
\item Supply actual timestamps or a verified timing-error bound, including
      channel synchronization and exposure averaging.
\item Establish stationary preparation and an independently justified
      dependence bound, or collect genuinely independent reset trials with
      their preparation cost included. Fitting the observed correlation does
      not establish hidden-state contraction.
\item Specify the physical state, time-reversal parity and relevant dynamical
      approximation. For a driven comparison, separate particle-bath entropy
      from the unrecorded energy of the detector, laser and controller.
\item Freeze the calibrated lag menu and statistical procedure before the
      next evaluation record. Then apply the dependence-aware finite-record
      confidence procedure to the bounded outputs.
\end{enumerate}

These requirements specify measurements that would make the theoretical
certificate usable. The present reanalysis already prevents two concrete
errors: treating a measured reference as exact latent position, and treating
small empirical moment departures as thermodynamic discoveries.
